%==========================================================================
\section{Conclusions}\label{sec:conclusions}

We redesigned an HBM statevector simulator for the Alveo U55C around streaming engines, on-chip tiles and a single multiplexed $2\times2$ primitive, and evaluated it with measured timing, power and fidelity on every platform. One card simulates 30 qubits in single precision and passes all \NumSuiteTotal{} cases of a functional suite started from entangled states. The redesign removes the single-beat memory access of the baseline kernel: end-to-end time at 24 qubits falls by \NumSpeedupVsevenMin to \NumSpeedupVsevenMax$\times$ and energy per circuit by \NumEnergyVsevenRatioMin to \NumEnergyVsevenRatioMax$\times$.

The measurements locate the design on the card's limits. One pseudo-channel sustains \NumHbmSixteenPerCURead~GB/s, the card \NumHbmThirtytwoCopy~GB/s over 32 channels, and the simulator's kernel time scales with the number of engines because each engine owns a channel. Circuits with few gates per pass are bound by that bandwidth, and for them one card is faster than a 128-core CPU server from \NumXplatGhzCrossN{} qubits. Circuits with many gates per pass are bound by engine arithmetic, and for them the server remains faster at 30 qubits. A data-centre GPU, the A100, is faster than the card by \NumXplatBestGpuRatioRqcThirty~to \NumXplatBestGpuRatioQftThirty$\times$ for random circuits and QFT at 30 qubits, about as fast for bandwidth-bound GHZ, and holds up to \NumGpuCudaqSingleMaxN~qubits.

The next steps follow from the same measurements: more arithmetic per engine or more engines to raise the compute bound, a host path in compiled code to reduce the millisecond overhead that dominates small circuits, and the distribution of the state over several cards, which is the subject of separate work~\cite{ali2026multifpga}.
